\chapter{ВИСНОВКИ}

У рамках дипломної роботи було розроблено платформу IoT MLOps для моніторингу промислових сенсорів у реальному часі з виявленням аномалій на основі машинного навчання.

Актуальність дослідження обумовлена потребою промислових підприємств у доступних рішеннях для моніторингу обладнання. Існуючі корпоративні платформи мають високу вартість та складність впровадження, тоді як рішення з відкритим кодом потребують значних зусиль для інтеграції компонентів у єдину систему. Це створює бар'єр для малих та середніх підприємств, які потребують надійного моніторингу, але не мають ресурсів для впровадження дорогих комерційних рішень.

Метою роботи було створення багатокористувацької платформи, що поєднує потокову обробку даних, машинне навчання для детекції аномалій та MLOps практики на базі технологій з відкритим кодом. Всі поставлені завдання було успішно виконано.

У ході роботи застосовано комплекс методів: аналітичний метод для вивчення існуючих IoT платформ та формулювання вимог; методи машинного навчання (Isolation Forest, Random Forest, XGBoost) для детекції аномалій у сенсорних даних; експериментальний метод для тестування продуктивності системи; порівняльний метод для оцінки алгоритмів та архітектурних рішень.

Основні результати роботи:

\begin{itemize}
    \item Розроблено мікросервісну архітектуру на базі контейнерів Docker, що забезпечує модульність та масштабованість системи
    \item Реалізовано потокову обробку даних з низькою затримкою через Apache Kafka та Apache Spark Structured Streaming
    \item Впроваджено архітектуру схема-на-тенанта для ізоляції даних між клієнтами з підтримкою стиснення TimescaleDB
    \item Створено MLOps інфраструктуру на базі MLflow для відстеження експериментів та версіонування моделей
    \item Досягнуто високу якість детекції аномалій на наборі даних Tennessee Eastman Process
    \item Розроблено систему сповіщень з підтримкою порогових правил та ML-детекції
    \item Створено веб-інтерфейс для візуалізації даних у реальному часі
\end{itemize}

Практична значущість роботи полягає у створенні готового рішення для промислових підприємств. Платформа забезпечує раннє виявлення несправностей обладнання, що дозволяє перейти від реактивного до превентивного обслуговування та знизити витрати на незаплановані простої. Використання технологій з відкритим кодом робить рішення доступним для підприємств з обмеженим бюджетом.

Наукова новизна полягає у поєднанні архітектури схема-на-тенанта з колонковим стисненням TimescaleDB, що забезпечує одночасно ізоляцію даних та ефективне використання сховища. Також запропоновано гібридний підхід до детекції аномалій, що поєднує порогові правила та моделі машинного навчання в єдиному механізмі сповіщень.

Напрямки подальшого розвитку включають: впровадження оркестрації Kubernetes для автоматичного масштабування; додавання моделей глибокого навчання (LSTM, Autoencoder) для детекції складних аномалій; розширення інтеграції з промисловими протоколами (OPC-UA, Modbus); реалізацію автоматичного перенавчання моделей при виявленні дрейфу даних.

Результати роботи підтверджують, що технології з відкритим кодом дозволяють створити конкурентоспроможну альтернативу комерційним IoT платформам для промислового моніторингу.
